% --- Archon physics formalization source begin ---
% archon:physics
% archon:covers PhyXMiniProblems/problem_phyx_mini_0968.lean
% archon:source-report reports/phyx_mini/problem_phyx_mini_0968.source.json

\chapter{Physics problem phyx\_mini\_0968}
\label{ch:PhyXMiniProblems_problem_phyx_mini_0968}

\paragraph{Problem source.}
\#\# Physical scenario

A pair of long, rigid metal rods, each of length 0.50 m, lie parallel to each other on a frictionless table. Their ends are connected by identical, very lightweight conducting springs with unstretched length \textbackslash{}( l\_0 \textbackslash{}) and force constant \textbackslash{}( k \textbackslash{}). When a current \textbackslash{}( I \textbackslash{}) runs through the circuit consisting of the rods and springs, the springs stretch. You measure the distance \textbackslash{}( x \textbackslash{}) each spring stretches for certain values of \textbackslash{}( I \textbackslash{}). When \textbackslash{}( I = 8.05\textbackslash{},\textbackslash{}text\{A\} \textbackslash{}), you measure that \textbackslash{}( x = 0.40\textbackslash{},\textbackslash{}text\{cm\} \textbackslash{}). When \textbackslash{}( I = 13.1\textbackslash{},\textbackslash{}text\{A\} \textbackslash{}), you find \textbackslash{}( x = 0.80\textbackslash{},\textbackslash{}text\{cm\} \textbackslash{}).

\#\# Question

What current is required for each spring to stretch \textbackslash{}( 1.00\textbackslash{},\textbackslash{}text\{cm\} \textbackslash{})?

\#\# Answer choices

- A: A: 12.0A

- B: B: 16.0A

- C: C: 13.1A

- D: D: 18.2A

\#\# Figure caption (auxiliary; use the image as primary evidence)

The image depicts two long, parallel wires with current flowing through them, indicated by the arrows and the symbol "I." The top arrow points to the right, while the bottom arrow points to the left. The wires are connected to springs on either side, suggesting they are movable or can exert force. The currents in the wires are flowing in opposite directions.

\#\# Dataset metadata

Magnetism; Multi-Formula Reasoning, Physical Model Grounding Reasoning

\paragraph{Recorded answer/context.}
B: B: 16.0A

\paragraph{Figure/image path.}
/root/proposal\_for\_physic/science-mango/phyx\_mini\_run/phyx\_data/test\_image/968.png

\paragraph{Formalization target.}
create a compiling Lean file with sorry bodies at `PhyXMiniProblems/problem\_phyx\_mini\_0968.lean`. The Lean declarations must preserve the physical quantities, dimensions or dimensional roles, figure labels, governing-law hypotheses, and final relation expressed by this problem.
Use Mathlib/Physlib names found through LeanExplore where available. If a physics API is missing, introduce faithful local abstractions rather than scalar placeholder aliases.

\begin{theorem}[Physics formalization target]
\label{thm:physics:phyx_mini_0968:target}
\lean{PhyXMiniProblems.ProblemPhyXMini0968.problem_phyx_mini_0968}
\uses{def:physics:phyx-mini-0968:phyxminiproblems-problemphyxmini0968-currentinamperes, def:physics:phyx-mini-0968:phyxminiproblems-problemphyxmini0968-parallelrodspringsetup, def:physics:phyx-mini-0968:phyxminiproblems-problemphyxmini0968-matcheswrittenparallelrodspringscenario, def:physics:phyx-mini-0968:phyxminiproblems-problemphyxmini0968-matchesprimaryparallelrodspringfigure, def:physics:phyx-mini-0968:phyxminiproblems-problemphyxmini0968-matchesparallelrodspringmeasurements, def:physics:phyx-mini-0968:phyxminiproblems-problemphyxmini0968-hasphysicalparallelrodspringparameters, def:physics:phyx-mini-0968:phyxminiproblems-problemphyxmini0968-satisfiesparallelrodspringlaws, def:physics:phyx-mini-0968:phyxminiproblems-problemphyxmini0968-recordeddatasetanswer, def:physics:phyx-mini-0968:phyxminiproblems-problemphyxmini0968-isuniqueclosestrequiredcurrentchoice}
The assigned autoformalize agent should translate this physics problem into Lean declarations in the covered file, with theorem and lemma proof bodies written as `by sorry`.
\end{theorem}
\begin{proof}
This is an autoformalization task, not a proof task. Produce statements that can later be proved by the physics prover without weakening the physical model.
\end{proof}

% --- Archon declaration topology begin ---
\section{Declaration topology}
\begin{definition}[electric Current Dimension]
  \label{def:physics:phyx-mini-0968:phyxminiproblems-problemphyxmini0968-electriccurrentdimension}
  \lean{PhyXMiniProblems.ProblemPhyXMini0968.electricCurrentDimension}
  Electric current has physical dimension charge per time
\end{definition}
\begin{proof}
  This is the declared model definition of electric Current Dimension.
\end{proof}
\begin{definition}[force Dimension]
  \label{def:physics:phyx-mini-0968:phyxminiproblems-problemphyxmini0968-forcedimension}
  \lean{PhyXMiniProblems.ProblemPhyXMini0968.forceDimension}
  Force has physical dimension M L T⁻²
\end{definition}
\begin{proof}
  This is the declared model definition of force Dimension.
\end{proof}
\begin{definition}[spring Stiffness Dimension]
  \label{def:physics:phyx-mini-0968:phyxminiproblems-problemphyxmini0968-springstiffnessdimension}
  \lean{PhyXMiniProblems.ProblemPhyXMini0968.springStiffnessDimension}
  A spring force constant has dimension force per length, namely M T⁻²
\end{definition}
\begin{proof}
  This is the declared model definition of spring Stiffness Dimension.
\end{proof}
\begin{definition}[magnetic Permeability Dimension]
  \label{def:physics:phyx-mini-0968:phyxminiproblems-problemphyxmini0968-magneticpermeabilitydimension}
  \lean{PhyXMiniProblems.ProblemPhyXMini0968.magneticPermeabilityDimension}
  Vacuum permeability has SI dimension N A⁻² = M L C⁻²
\end{definition}
\begin{proof}
  This is the declared model definition of magnetic Permeability Dimension.
\end{proof}
\begin{definition}[Length Magnitude]
  \label{def:physics:phyx-mini-0968:phyxminiproblems-problemphyxmini0968-lengthmagnitude}
  \lean{PhyXMiniProblems.ProblemPhyXMini0968.LengthMagnitude}
  A nonnegative, unit-independent physical length
\end{definition}
\begin{proof}
  This is the declared model definition of Length Magnitude.
\end{proof}
\begin{definition}[Electric Current Magnitude]
  \label{def:physics:phyx-mini-0968:phyxminiproblems-problemphyxmini0968-electriccurrentmagnitude}
  \lean{PhyXMiniProblems.ProblemPhyXMini0968.ElectricCurrentMagnitude}
  \uses{def:physics:phyx-mini-0968:phyxminiproblems-problemphyxmini0968-electriccurrentdimension}
  A nonnegative, unit-independent electric-current magnitude
\end{definition}
\begin{proof}
  This is the declared model definition of Electric Current Magnitude.
\end{proof}
\begin{definition}[Force Magnitude]
  \label{def:physics:phyx-mini-0968:phyxminiproblems-problemphyxmini0968-forcemagnitude}
  \lean{PhyXMiniProblems.ProblemPhyXMini0968.ForceMagnitude}
  \uses{def:physics:phyx-mini-0968:phyxminiproblems-problemphyxmini0968-forcedimension}
  A nonnegative, unit-independent force magnitude
\end{definition}
\begin{proof}
  This is the declared model definition of Force Magnitude.
\end{proof}
\begin{definition}[Spring Stiffness Magnitude]
  \label{def:physics:phyx-mini-0968:phyxminiproblems-problemphyxmini0968-springstiffnessmagnitude}
  \lean{PhyXMiniProblems.ProblemPhyXMini0968.SpringStiffnessMagnitude}
  \uses{def:physics:phyx-mini-0968:phyxminiproblems-problemphyxmini0968-springstiffnessdimension}
  A nonnegative, unit-independent spring stiffness
\end{definition}
\begin{proof}
  This is the declared model definition of Spring Stiffness Magnitude.
\end{proof}
\begin{definition}[Magnetic Permeability Magnitude]
  \label{def:physics:phyx-mini-0968:phyxminiproblems-problemphyxmini0968-magneticpermeabilitymagnitude}
  \lean{PhyXMiniProblems.ProblemPhyXMini0968.MagneticPermeabilityMagnitude}
  \uses{def:physics:phyx-mini-0968:phyxminiproblems-problemphyxmini0968-magneticpermeabilitydimension}
  A nonnegative, unit-independent magnetic permeability
\end{definition}
\begin{proof}
  This is the declared model definition of Magnetic Permeability Magnitude.
\end{proof}
\begin{definition}[coherent SIReadout]
  \label{def:physics:phyx-mini-0968:phyxminiproblems-problemphyxmini0968-coherentsireadout}
  \lean{PhyXMiniProblems.ProblemPhyXMini0968.coherentSIReadout}
  Read a nonnegative dimensionful quantity in coherent SI units
\end{definition}
\begin{proof}
  This is the declared model definition of coherent SIReadout.
\end{proof}
\begin{definition}[length Readout]
  \label{def:physics:phyx-mini-0968:phyxminiproblems-problemphyxmini0968-lengthreadout}
  \lean{PhyXMiniProblems.ProblemPhyXMini0968.lengthReadout}
  \uses{def:physics:phyx-mini-0968:phyxminiproblems-problemphyxmini0968-lengthmagnitude}
  Read a physical length in a selected length unit
\end{definition}
\begin{proof}
  This is the declared model definition of length Readout.
\end{proof}
\begin{definition}[length In Meters]
  \label{def:physics:phyx-mini-0968:phyxminiproblems-problemphyxmini0968-lengthinmeters}
  \lean{PhyXMiniProblems.ProblemPhyXMini0968.lengthInMeters}
  \uses{def:physics:phyx-mini-0968:phyxminiproblems-problemphyxmini0968-lengthmagnitude, def:physics:phyx-mini-0968:phyxminiproblems-problemphyxmini0968-coherentsireadout}
  Coherent-SI length readout, in metres
\end{definition}
\begin{proof}
  This is the declared model definition of length In Meters.
\end{proof}
\begin{definition}[current In Amperes]
  \label{def:physics:phyx-mini-0968:phyxminiproblems-problemphyxmini0968-currentinamperes}
  \lean{PhyXMiniProblems.ProblemPhyXMini0968.currentInAmperes}
  \uses{def:physics:phyx-mini-0968:phyxminiproblems-problemphyxmini0968-electriccurrentmagnitude, def:physics:phyx-mini-0968:phyxminiproblems-problemphyxmini0968-coherentsireadout}
  Coherent-SI current readout, in amperes
\end{definition}
\begin{proof}
  This is the declared model definition of current In Amperes.
\end{proof}
\begin{definition}[force In Newtons]
  \label{def:physics:phyx-mini-0968:phyxminiproblems-problemphyxmini0968-forceinnewtons}
  \lean{PhyXMiniProblems.ProblemPhyXMini0968.forceInNewtons}
  \uses{def:physics:phyx-mini-0968:phyxminiproblems-problemphyxmini0968-forcemagnitude, def:physics:phyx-mini-0968:phyxminiproblems-problemphyxmini0968-coherentsireadout}
  Coherent-SI force readout, in newtons
\end{definition}
\begin{proof}
  This is the declared model definition of force In Newtons.
\end{proof}
\begin{definition}[spring Stiffness In Newtons Per Meter]
  \label{def:physics:phyx-mini-0968:phyxminiproblems-problemphyxmini0968-springstiffnessinnewtonspermeter}
  \lean{PhyXMiniProblems.ProblemPhyXMini0968.springStiffnessInNewtonsPerMeter}
  \uses{def:physics:phyx-mini-0968:phyxminiproblems-problemphyxmini0968-springstiffnessmagnitude, def:physics:phyx-mini-0968:phyxminiproblems-problemphyxmini0968-coherentsireadout}
  Coherent-SI spring-stiffness readout, in newtons per metre
\end{definition}
\begin{proof}
  This is the declared model definition of spring Stiffness In Newtons Per Meter.
\end{proof}
\begin{definition}[magnetic Permeability In Newtons Per Ampere Squared]
  \label{def:physics:phyx-mini-0968:phyxminiproblems-problemphyxmini0968-magneticpermeabilityinnewtonsperamperesquared}
  \lean{PhyXMiniProblems.ProblemPhyXMini0968.magneticPermeabilityInNewtonsPerAmpereSquared}
  \uses{def:physics:phyx-mini-0968:phyxminiproblems-problemphyxmini0968-magneticpermeabilitymagnitude, def:physics:phyx-mini-0968:phyxminiproblems-problemphyxmini0968-coherentsireadout}
  Coherent-SI vacuum-permeability readout, in newtons per ampere squared
\end{definition}
\begin{proof}
  This is the declared model definition of magnetic Permeability In Newtons Per Ampere Squared.
\end{proof}
\begin{definition}[Rod Label]
  \label{def:physics:phyx-mini-0968:phyxminiproblems-problemphyxmini0968-rodlabel}
  \lean{PhyXMiniProblems.ProblemPhyXMini0968.RodLabel}
  The upper and lower rods distinguished in the supplied image
\end{definition}
\begin{proof}
  This is the declared model definition of Rod Label.
\end{proof}
\begin{definition}[Spring Label]
  \label{def:physics:phyx-mini-0968:phyxminiproblems-problemphyxmini0968-springlabel}
  \lean{PhyXMiniProblems.ProblemPhyXMini0968.SpringLabel}
  The two identical springs, named by their position in the image
\end{definition}
\begin{proof}
  This is the declared model definition of Spring Label.
\end{proof}
\begin{definition}[Rod End]
  \label{def:physics:phyx-mini-0968:phyxminiproblems-problemphyxmini0968-rodend}
  \lean{PhyXMiniProblems.ProblemPhyXMini0968.RodEnd}
  The two ends of each horizontal rod
\end{definition}
\begin{proof}
  This is the declared model definition of Rod End.
\end{proof}
\begin{definition}[Horizontal Direction]
  \label{def:physics:phyx-mini-0968:phyxminiproblems-problemphyxmini0968-horizontaldirection}
  \lean{PhyXMiniProblems.ProblemPhyXMini0968.HorizontalDirection}
  Horizontal directions used by the current arrows
\end{definition}
\begin{proof}
  This is the declared model definition of Horizontal Direction.
\end{proof}
\begin{definition}[Operating Point]
  \label{def:physics:phyx-mini-0968:phyxminiproblems-problemphyxmini0968-operatingpoint}
  \lean{PhyXMiniProblems.ProblemPhyXMini0968.OperatingPoint}
  The three operating points named in the problem statement
\end{definition}
\begin{proof}
  This is the declared model definition of Operating Point.
\end{proof}
\begin{definition}[Parallel Rod Spring Figure]
  \label{def:physics:phyx-mini-0968:phyxminiproblems-problemphyxmini0968-parallelrodspringfigure}
  \lean{PhyXMiniProblems.ProblemPhyXMini0968.ParallelRodSpringFigure}
  \uses{def:physics:phyx-mini-0968:phyxminiproblems-problemphyxmini0968-rodlabel, def:physics:phyx-mini-0968:phyxminiproblems-problemphyxmini0968-springlabel, def:physics:phyx-mini-0968:phyxminiproblems-problemphyxmini0968-rodend, def:physics:phyx-mini-0968:phyxminiproblems-problemphyxmini0968-horizontaldirection}
  Literal qualitative content of image 968.png. It shows two horizontal parallel rods, a spring joining their left ends and another joining their right ends, and purple current arrows labeled I in opposite directions
\end{definition}
\begin{proof}
  This is the declared model definition of Parallel Rod Spring Figure.
\end{proof}
\begin{definition}[Parallel Rod Spring Setup]
  \label{def:physics:phyx-mini-0968:phyxminiproblems-problemphyxmini0968-parallelrodspringsetup}
  \lean{PhyXMiniProblems.ProblemPhyXMini0968.ParallelRodSpringSetup}
  \uses{def:physics:phyx-mini-0968:phyxminiproblems-problemphyxmini0968-lengthmagnitude, def:physics:phyx-mini-0968:phyxminiproblems-problemphyxmini0968-electriccurrentmagnitude, def:physics:phyx-mini-0968:phyxminiproblems-problemphyxmini0968-forcemagnitude, def:physics:phyx-mini-0968:phyxminiproblems-problemphyxmini0968-springstiffnessmagnitude, def:physics:phyx-mini-0968:phyxminiproblems-problemphyxmini0968-magneticpermeabilitymagnitude, def:physics:phyx-mini-0968:phyxminiproblems-problemphyxmini0968-operatingpoint, def:physics:phyx-mini-0968:phyxminiproblems-problemphyxmini0968-parallelrodspringfigure}
  Independent physical quantities and force observables for the apparatus. In particular, the current at the requested extension is not defined from an answer choice. The force observables are constrained only by the governing laws below
\end{definition}
\begin{proof}
  This is the declared model definition of Parallel Rod Spring Setup.
\end{proof}
\begin{definition}[Matches Written Parallel Rod Spring Scenario]
  \label{def:physics:phyx-mini-0968:phyxminiproblems-problemphyxmini0968-matcheswrittenparallelrodspringscenario}
  \lean{PhyXMiniProblems.ProblemPhyXMini0968.MatchesWrittenParallelRodSpringScenario}
  \uses{def:physics:phyx-mini-0968:phyxminiproblems-problemphyxmini0968-parallelrodspringsetup}
  Qualitative properties stated in the written physical scenario
\end{definition}
\begin{proof}
  This is the declared model definition of Matches Written Parallel Rod Spring Scenario.
\end{proof}
\begin{definition}[Matches Primary Parallel Rod Spring Figure]
  \label{def:physics:phyx-mini-0968:phyxminiproblems-problemphyxmini0968-matchesprimaryparallelrodspringfigure}
  \lean{PhyXMiniProblems.ProblemPhyXMini0968.MatchesPrimaryParallelRodSpringFigure}
  \uses{def:physics:phyx-mini-0968:phyxminiproblems-problemphyxmini0968-parallelrodspringsetup}
  Primary-raster evidence only. These fields record object placement and arrow directions; they contain no numerical current or target-answer information
\end{definition}
\begin{proof}
  This is the declared model definition of Matches Primary Parallel Rod Spring Figure.
\end{proof}
\begin{definition}[Matches Parallel Rod Spring Measurements]
  \label{def:physics:phyx-mini-0968:phyxminiproblems-problemphyxmini0968-matchesparallelrodspringmeasurements}
  \lean{PhyXMiniProblems.ProblemPhyXMini0968.MatchesParallelRodSpringMeasurements}
  \uses{def:physics:phyx-mini-0968:phyxminiproblems-problemphyxmini0968-lengthreadout, def:physics:phyx-mini-0968:phyxminiproblems-problemphyxmini0968-lengthinmeters, def:physics:phyx-mini-0968:phyxminiproblems-problemphyxmini0968-currentinamperes, def:physics:phyx-mini-0968:phyxminiproblems-problemphyxmini0968-parallelrodspringsetup}
  Numerical problem data. The two measured currents and extensions calibrate the unknown spring and separation parameters. The requested operating point contains only its prescribed extension; its current is deliberately absent
\end{definition}
\begin{proof}
  This is the declared model definition of Matches Parallel Rod Spring Measurements.
\end{proof}
\begin{definition}[Has Physical Parallel Rod Spring Parameters]
  \label{def:physics:phyx-mini-0968:phyxminiproblems-problemphyxmini0968-hasphysicalparallelrodspringparameters}
  \lean{PhyXMiniProblems.ProblemPhyXMini0968.HasPhysicalParallelRodSpringParameters}
  \uses{def:physics:phyx-mini-0968:phyxminiproblems-problemphyxmini0968-lengthinmeters, def:physics:phyx-mini-0968:phyxminiproblems-problemphyxmini0968-springstiffnessinnewtonspermeter, def:physics:phyx-mini-0968:phyxminiproblems-problemphyxmini0968-magneticpermeabilityinnewtonsperamperesquared, def:physics:phyx-mini-0968:phyxminiproblems-problemphyxmini0968-parallelrodspringsetup}
  Positivity and nondegeneracy conditions for the physical apparatus
\end{definition}
\begin{proof}
  This is the declared model definition of Has Physical Parallel Rod Spring Parameters.
\end{proof}
\begin{definition}[Satisfies Parallel Rod Spring Laws]
  \label{def:physics:phyx-mini-0968:phyxminiproblems-problemphyxmini0968-satisfiesparallelrodspringlaws}
  \lean{PhyXMiniProblems.ProblemPhyXMini0968.SatisfiesParallelRodSpringLaws}
  \uses{def:physics:phyx-mini-0968:phyxminiproblems-problemphyxmini0968-lengthinmeters, def:physics:phyx-mini-0968:phyxminiproblems-problemphyxmini0968-currentinamperes, def:physics:phyx-mini-0968:phyxminiproblems-problemphyxmini0968-forceinnewtons, def:physics:phyx-mini-0968:phyxminiproblems-problemphyxmini0968-springstiffnessinnewtonspermeter, def:physics:phyx-mini-0968:phyxminiproblems-problemphyxmini0968-magneticpermeabilityinnewtonsperamperesquared, def:physics:phyx-mini-0968:phyxminiproblems-problemphyxmini0968-parallelrodspringsetup}
  School-level geometry, Hooke, long-parallel-wire, and equilibrium laws. The magnetic law gives the magnitude on either rod when the equal currents are oppositely directed; hence the force is repulsive as shown by the stretching springs. These are uniform laws over all operating points, not the solved formula for the requested current
\end{definition}
\begin{proof}
  This is the declared model definition of Satisfies Parallel Rod Spring Laws.
\end{proof}
\begin{definition}[Answer Choice]
  \label{def:physics:phyx-mini-0968:phyxminiproblems-problemphyxmini0968-answerchoice}
  \lean{PhyXMiniProblems.ProblemPhyXMini0968.AnswerChoice}
  Labels of the four current choices printed with the problem
\end{definition}
\begin{proof}
  This is the declared model definition of Answer Choice.
\end{proof}
\begin{definition}[answer Current In Amperes]
  \label{def:physics:phyx-mini-0968:phyxminiproblems-problemphyxmini0968-answercurrentinamperes}
  \lean{PhyXMiniProblems.ProblemPhyXMini0968.answerCurrentInAmperes}
  \uses{def:physics:phyx-mini-0968:phyxminiproblems-problemphyxmini0968-answerchoice}
  Current in amperes printed beside each answer label
\end{definition}
\begin{proof}
  This is the declared model definition of answer Current In Amperes.
\end{proof}
\begin{definition}[recorded Dataset Answer]
  \label{def:physics:phyx-mini-0968:phyxminiproblems-problemphyxmini0968-recordeddatasetanswer}
  \lean{PhyXMiniProblems.ProblemPhyXMini0968.recordedDatasetAnswer}
  \uses{def:physics:phyx-mini-0968:phyxminiproblems-problemphyxmini0968-answerchoice}
  The answer label recorded in the source dataset
\end{definition}
\begin{proof}
  This is the declared model definition of recorded Dataset Answer.
\end{proof}
\begin{definition}[Is Unique Closest Required Current Choice]
  \label{def:physics:phyx-mini-0968:phyxminiproblems-problemphyxmini0968-isuniqueclosestrequiredcurrentchoice}
  \lean{PhyXMiniProblems.ProblemPhyXMini0968.IsUniqueClosestRequiredCurrentChoice}
  \uses{def:physics:phyx-mini-0968:phyxminiproblems-problemphyxmini0968-currentinamperes, def:physics:phyx-mini-0968:phyxminiproblems-problemphyxmini0968-parallelrodspringsetup, def:physics:phyx-mini-0968:phyxminiproblems-problemphyxmini0968-answerchoice, def:physics:phyx-mini-0968:phyxminiproblems-problemphyxmini0968-answercurrentinamperes}
  A choice is strictly nearer to the required current than every rival
\end{definition}
\begin{proof}
  This is the declared model definition of Is Unique Closest Required Current Choice.
\end{proof}
\begin{lemma}[calibrated Length And Required Current Squared]
  \label{lem:physics:phyx-mini-0968:phyxminiproblems-problemphyxmini0968-calibratedlengthandrequiredcurrentsquared}
  \lean{PhyXMiniProblems.ProblemPhyXMini0968.calibratedLengthAndRequiredCurrentSquared}
  \uses{def:physics:phyx-mini-0968:phyxminiproblems-problemphyxmini0968-lengthinmeters, def:physics:phyx-mini-0968:phyxminiproblems-problemphyxmini0968-currentinamperes, def:physics:phyx-mini-0968:phyxminiproblems-problemphyxmini0968-parallelrodspringsetup, def:physics:phyx-mini-0968:phyxminiproblems-problemphyxmini0968-matchesparallelrodspringmeasurements, def:physics:phyx-mini-0968:phyxminiproblems-problemphyxmini0968-hasphysicalparallelrodspringparameters, def:physics:phyx-mini-0968:phyxminiproblems-problemphyxmini0968-satisfiesparallelrodspringlaws}
  Eliminating the common spring stiffness and magnetic coefficient from the two calibration measurements gives l₀ = 1752/210025 m. Applying the same laws at x = 1.00 cm gives I² = 15409/64 A². Both are conclusions rather than input data
\end{lemma}
\begin{proof}
  The conclusion follows from the governing relations and figure readouts named in the statement of calibrated Length And Required Current Squared.
\end{proof}
% --- Archon declaration topology end ---
% --- Archon physics formalization source end ---
